\section{Operating problem and methods}

\subsection{Service concept}
The service is a bidirectional, headway-managed route for passengers who arrive without targeting a published departure or a timed transfer. The original nominal roster provides 131 trips per direction at a 360-second nominal spacing. Vehicles come from a single fixed physical fleet, reversing direction at each terminal. Every 900 seconds the planner updates the unreleased departure plan; a trip already released, queried for dispatch, or eligible for a bus is not moved. The projection assigns an active block of gaps and a compensating tail while retaining the trip roster and final departure no earlier than the demand cutoff. Dispatch also requires a ready vehicle and a 120-second minimum actual spacing. Neither that floor nor the planned spacing guarantees a particular realized headway distribution.

The model commits trips at the simulator decision boundary. It does not promise a public notice period and does not model a passenger reaction to announcements. In operation, a headway-managed route would need rolling passenger information and an agency-defined lock for already advertised departures. For timetable-based service or guaranteed connections, early departures would need a separate constraint: a plan that runs with little readiness delay may still be unsuitable. The untimed-arrival assumption is most plausible in frequent periods. The study does not assume it remains valid during low-frequency shoulders or near last departures; demand generation stops at 19:03, and the roster and evaluation include a separate clearance period.

\subsection{Controller and observations}
The upper policy selects among 15 actions every 15 minutes. Each is a pair of a quota-adjusted spacing multiplier in $\{0.78,0.89,1.00,1.12,1.24\}$ and a directional tilt in $\{-36,0,36\}$ seconds. The inventory guard attenuates the tilt when directional ready stock is low. The same projection is used for all corrected comparisons. Along the route, decisions are made after boarding and alighting using the state of the bus and causal estimates of the adjacent gaps. Stop holding is restricted to the bins $\{0,5,10,15,20,25\}$ seconds.

The turnaround ledger predicts ready times using current terminal stock, active vehicle position, residual stop time, and a recent direction-specific median nonholding trip time. At dispatch and at least every 60 seconds, it sorts predictions at each terminal and matches them to unlaunched trips. It deducts a 20-second planning buffer from predicted slack and converts the positive remainder to a budget with retention 0.50. The budget is capped at 85 seconds; stronger opposite-direction demand and less than 75 seconds of slack instead trigger a 0.40 multiplier on the allowance and a 30-second cap. An initially unmatched trip gets 37.5 seconds. A later calculation can shrink but not replenish the budget, and actual holding is debited. This is a hard action limit, not a guarantee that a predicted vehicle will return on time.

The causal boarding proxy divides observed boardings by the observed arrival gap (with fallback, clipping, and a fullness adjustment) and updates a held stop rate. A 20-second sample contributes one third of its per-minute rate to the closed one-minute bin. A four-harmonic recursive least-squares estimator with a static demand prior separates the smooth component and residual. The aligned policy sends directional low-frequency summaries to the upper actor and signed local residual summaries to the lower actor. Both use shared physical features. The upper state contains 13 physical and clock features and eight demand features; the lower state has 15 physical and clock features, including the mask limit, and four demand features. Neither actor reads the station's exact post-boarding queue. Swapped uses matched summaries of the opposite band. NoFreq uses the same summaries of the completed, unfiltered APC bin and its one-bin change for both actors. The ledger still reads the same harmonic low-frequency signals in these variants; the comparisons change actor inputs, not all harmonic information in the system.

Discrete entropy-regularized actor--critic policies \citep{haarnoja2018soft,christodoulou2019discrete} are trained with the same critic and update budget. The upper reward penalizes simulated waiting, exposure to discretionary holding, an excess-waiting proxy, and actual vehicle-readiness delay; it also includes a remaining-quota potential. The lower reward subtracts onboard delay and predicted slack penalties from a two-gap waiting proxy. With observed forward and predicted follower gaps $g_f,g_b$, targets $H_f,H_b$, and a downstream boarding-rate proxy $\lambda$, the two-sided term is
\begin{equation}
G(u)=\lambda\frac{H_f+H_b}{2}u
  \left[\frac{g_b}{H_b}-\frac{g_f}{H_f}
    -\frac{u}{2}\left(\frac1{H_f}+\frac1{H_b}\right)\right].
\end{equation}
At training time the downstream rates are reconstructed from subsequent boardings on that trip. This changes a reward label, not a runtime observation. The lower critic uses a zero-hold action reference and a Lagrange penalty for regularity overshoot, onboard load, and tight projected slack; neither the proxy nor the cost critic ensures a non-harmful action in operation.

\subsection{Benchmarks and evaluation}
The corrected case uses Fall 2025 seasonal automatic passenger count averages and stop geometry for MBTA Route 111. It is a 23-stop data-informed simulation: the published trips are not replayed, origin--destination flows are reconstructed, running speeds are constructed, and control uses a synthetic timetable. Exogenous traffic and demand are paired by scenario tape. V503 swaps the actor frequency inputs on 16 held-out days. V504 evaluates four 200-episode seeds over a different 16-day test set against both matching-budget, matching-architecture NoFreq and a dynamic rule (Manual). Manual picks from the same 15 release actions by nearest directional square-root target based on the raw completed boarding bin, and uses a causal, load-aware hold scorer based on the same waiting and slack proxy. It runs with the same six holding bins, projection, and budget. Shared constraints do not imply matched actual holding dose.

To attribute holding separately, a closed-loop intervention keeps each Aligned upper actor fixed and substitutes either zero discretionary holding or the Manual hold scorer for its lower policy. The upper actor replans from that arm's own realized vehicle stock and holding history. Simply replaying an identical departure plan would suppress this feedback. The intervention reuses the same 16 V504 days, so it is a mechanism comparison, not a new external confirmation. The additional Rolling baseline combines the same Manual hold scorer with an upper search over the native 15 projected plans. It forecasts the next eight trips per direction, recycles the predicted returning fleet, and scores squared-gap expected waiting plus readiness cost (one nominal gap's arrivals per second of delay). Its inputs are the same causal stock, position, current stop time, recent nonholding times, and completed APC bins. It is not a reproduction of the integer nonlinear program of \citet{gkiotsalitis2021rescheduling}.

We report waiting plus onboard exposure per generated passenger, actual arrival-headway CV$^2$, and actual vehicle-readiness delay from first dispatch eligibility under the active plan. Exposure includes passengers still waiting or onboard at the fixed end of evaluation. We also report the absolute actual departure displacement from the original nominal timetable, decomposed into positive earliness and lateness, and note whether every scheduled trip and passenger was served. The common original timetable is an operating reference, not a passenger promise in this simulation. Paired bootstraps resample training seeds and demand days independently (30,000 draws), with deterministic rule scores shared rather than treated as additional trained seeds. The joint criterion requires the upper 95\% endpoint to be negative for journey and nonpositive for CV$^2$, readiness and absolute deviation (V503 uses journey, CV$^2$ and readiness).
